\documentclass[a4paper]{article}
\usepackage[utf8]{inputenc}
\usepackage{csquotes} %untuk mengubah tanda kutip(biblatex butuh ini)
\usepackage[backend=biber,style=verbose]{biblatex} %untuk bibiliography(daftar pustaka dan footnote)
%verbose karena style "apa" ada masalah di device saya
\usepackage{amsmath} %untuk simbol matematika dasar
\usepackage{lipsum} %untuk dummy text
\usepackage{ulem}%untuk underline
\usepackage{graphicx}
\usepackage{geometry}
\geometry{left=2.5cm, right=2.5cm, top=3cm, bottom=3cm}
\usepackage[indonesian]{babel}
\usepackage{xcolor} % untuk text color

%definisi warna warna
\definecolor{vermilion}{RGB}{227,66,52}
\definecolor{amber}{RGB}{255,192,0}
\definecolor{chartreuse}{RGB}{128,255,0}
\definecolor{teal}{RGB}{0,128,128}
\definecolor{violet}{RGB}{128,0,255}
\definecolor{russet}{RGB}{128,70,27}
\definecolor{slate}{RGB}{112,128,144}
\definecolor{citron}{RGB}{221,208,106}
\definecolor{scarlet}{RGB}{255,36,0}
\definecolor{spring-green}{RGB}{0,255,128}
\definecolor{azure}{RGB}{0,128,255}
\definecolor{rose}{RGB}{255,0,128}

\usepackage{hyperref} % Untuk hyperlink
\usepackage{tocloft} % Untuk mengatur daftar isi dari suatu daftar (bukan table of content)

\usepackage{titling}
\usepackage{tikz} % untuk grafis dan pembuatan sampel warna
\usetikzlibrary{calc}

% Untuk mengubah format penomoran menjadi arab
\renewcommand\thesection{\arabic{section}}
% Untuk mengubah format penomoran agar mengabaikan label dari format penomoran section
\renewcommand\thesubsection{\arabic{section}.\arabic{subsection}}

% Untuk daftar pustaka pemformatan nomor dihapus
\newcommand{\nonumsection}[1]{\section*{#1}
\addcontentsline{toc}{section}{#1}
}
\usepackage{xparse} % flexible new command serta logical IfValueTF
\usepackage{mfirstuc} % kapital huruf pertama
\addbibresource{src/makalah.bib} %file daftar pustaka

% Membuat textbox dengan background warna tertentu
\NewDocumentCommand{\warnaln}{mo}{%
  \tikz[baseline=-0.5\ht\strutbox]{%
    \node[shape=rectangle, draw=black, fill=#1, minimum height=1.5em, inner sep=2pt]{%
      \IfValueTF{#2}{\capitalisewords{#2}}{%
      \capitalisewords{#1}
      }%
    };
  }%
}

\newcommand{\warnalnbl}[1]{
  \tikz[baseline=-0.5\ht\strutbox]{
    \node[shape=rectangle, draw=black, fill=black, minimum height=1.5em, inner sep= 2pt, text=white]{#1}
  }
}

\title{%
  \vspace{-3cm}
  \centering
  \vspace{1cm}
  \textbf{Judul Makalah}\\
  \large Mata Kuliah: \textbf{Nama Mata Kuliah}\\
}
\author{$\begin{array}{c}
  \textbf{M Reza Dwi Prasetiawan}\\
  \textbf{Nama anggota lain}\\
  \textbf{Nama anggota lain}\\
  \textbf{Nama anggota lain}\\
\end{array}$
}
\date{\today}

\begin{document}
\pagenumbering{gobble}
\begin{titlepage}
  \maketitle
  \vfill
  \begin{center}
    \includegraphics[width=0.6\textwidth]{resources/logo-univ.png}\\
    \large
    Teknologi Informasi\\
    Nama Univ\\
    Lanjutan nama univ\\
    Dosen Pengampu: \textbf{Nama Dosen, gelar}\\
  \end{center}
\end{titlepage}

\newpage
\pagenumbering{arabic}
% Insert table of contents
\tableofcontents
\newpage
\section{Pendahuluan}
\lipsum[1]
\subsection{Latar Belakang}
\lipsum[2-3]\footcite{book1}
\subsection{Rumusan Masalah}
\lipsum[4]
\newpage

\section{Metode Penelitian}
\lipsum[5]\footcite{book2}
\subsection{Metode Penelitian 1}
\lipsum[6]\footcite{book3}
\subsection{Metode Penelitian 2}
\lipsum[7]\footcite{book4}
\subsection{Metode Penelitian 3}
\lipsum[8]\footcite{book5}
\newpage

\section{Kesimpulan}
\lipsum[4-5]\footcite{book1}
\footcite{book2}
\newpage

\nonumsection{Daftar Pustaka}
\printbibliography[heading=none, title={Daftar Pustaka}]
\begin{enumerate}
  \item \textcolor{blue}{\href{link}{Hint Text}}
\end{enumerate}
\end{document}

